\chapter{Koncepció}

\section{A vizsgált probléma kör}

A kimenetek valószínűsége és ahhoz tartozó kifizetések határozza meg, hogy milyen hosszútávú eredményt fog adni egy kaszinójáték, ez a matematika vizsgálatnak a legfontosabb kérdése.
Vannak különböző volatilitású játékok. Az egyik lehet sokkal gyakrabban ad kisebb nyereményeket. A másikban visszafizetések ritkábbak, de annál nagyobb összegűek. Emiatt a nyerések gyakorisága
nem elegendő az összehasonlításhoz

A korábban említett visszafizetési arány(RTP), az általunk megjátszott tétnek mekkora része fog megtérülni. Ez egy nagyon hosszú távú tulajdonság. A megjátszott tétek és a nyeremények összegének hányadosával kapjuk meg.
Egy rövid játéksorozat alkalmával nem valószínű, hogy visszakapjuk. Ez egy hosszútávú tulajdonság. Az ingadozás miatt vizsgálni kell a nyeremények eloszlását és szórását.

A célünk az alkalmazással a játék kipróbálásának lehetőségét összekapcsolja a fent említett összefüggését. A felhasználó saját pénz veszteség nélkül játszhat, a rendszer pedig eltárolja a körök költségét,
jóváírását és a nettó haszont. Az ebből keletkező adatok további vizsgálatokhoz is felhasználható.

\section{A játékok matematikai háttere}

A játékok eltérő valószínűségi modellekel működnek. Emiatt különböző számítási és elemzési módszerek bevonása szükséges.

Európai rulett kiosztást alkalmazunk, egyetlen eredményt választunk 0-tól 36-ig. ha minden szám valószínűsége azonos, akkor egyenletes eloszlás.
Egy fogadás nyerési esélye a megfogadott számokból, a várható nyeremény a találat esélyéből és a kifizetés kapcsolatából számítható ki.


\section{Amit csak említés szintjén érdemes szerepeltetni}

Az olvasóról annyit feltételezhetünk, hogy programozásban valamilyen szinten járatos, és a matematikai alapfogalmakkal sem ebben a dolgozatban kell megismertetni.
A speciális eszközök, programozási nyelvek, matematikai módszerek és jelölések persze jó, hogy ha említésre kerülnek, de nem kell nagyon belemenni a közismertnek tekinthető dolgokba.
